\documentclass[12pt]{article}
\usepackage[T1]{fontenc}
\usepackage[utf8]{inputenc}
\usepackage{lmodern}
\usepackage{textcomp}
\usepackage{enumitem}
\usepackage{geometry}
\geometry{margin=1in}
\begin{document}
\begin{center}
Answer Key - History - Graduate Level Assessment 11 \
\small (Answers are ordered exactly as in the question paper.)
\end{center}

\section*{Multiple Choice}
\begin{enumerate}[label=\arabic*.]
\item \textbf{Answer:} D
\\ \textit{Options:} A) meat to fruits and vegetables.; B) species focused on grasses to fruits and vegetables.; C) nuts and fruits to fish.; D) nuts and fruits to species more focused on grasses.; E) species focused on grasses to meat.
\\ \textit{Note:} The correct answer reflects the understanding that the shift from C$_{3}$ to C$_{4}$ photosynthetic pathways in plants, which are associated with warmer, drier climates, led to a decrease in forested areas and an increase in grasslands, thereby altering ancient hominid diets from reliance on nuts and fruits to an emphasis on grasses as a food source.
\item \textbf{Answer:} C
\\ \textit{Options:} A) instructions for the builders written there by the engineers.; B) random scratches and not a meaningful script or written language.; C) intended to identify the local group that had produced them.; D) simple arithmetic performed by carpenters on the building site.
\\ \textit{Note:} The markings on the mud-bricks of Huaca del Sol likely served as a means of identification to signify the specific local group responsible for their production, reflecting the social and cultural practices of the time.
\item \textbf{Answer:} C
\\ \textit{Options:} A) in sub-Saharan Africa, about 10,500 years ago; B) in western Europe, about 3,500 years ago; C) in the Middle East, about 10,500 years ago; D) in North America, about 5,500 years ago; E) in East Asia, about 5,500 years ago
\\ \textit{Note:} The correct answer is supported by archaeological findings that trace the origins of cattle domestication to the Fertile Crescent in the Middle East, specifically around 10,500 years ago, marking a significant development in human agriculture and society.
\item \textbf{Answer:} D
\\ \textit{Options:} A) trade.; B) warfare.; C) the development of a common religion.; D) conspicuous consumption by elites.
\\ \textit{Note:} The Minoan culture is characterized by a relatively egalitarian society with limited evidence of wealth disparity or luxurious displays of material goods typically associated with elite classes in other early civilizations, indicating a lack of conspicuous consumption.
\item \textbf{Answer:} E
\\ \textit{Options:} A) Erectus was not bipedal.; B) Erectus was not capable of using tools.; C) Erectus was primarily herbivorous.; D) Erectus fossils are found only in Africa.; E) Erectus possessed a larger brain.
\\ \textit{Note:} The correct answer is based on the fact that Homo erectus had a significantly larger brain size, averaging around 900 to 1,100 cubic centimeters, compared to Homo habilis, which had a brain size of approximately 510 to 600 cubic centimeters, indicating a higher level of cognitive development and more advanced tool-making capabilities.
\end{enumerate}

\section*{True / False}
\begin{enumerate}[label=\arabic*.]
\setcounter{enumi}{5}
\item \textbf{Answer:} False
\\ \textit{Options:} A) True; B) False
\\ \textit{Note:} The correct answer is: 2009
\item \textbf{Answer:} False
\\ \textit{Options:} A) True; B) False
\\ \textit{Note:} The correct answer is: Vistula River
\item \textbf{Answer:} False
\\ \textit{Options:} A) True; B) False
\\ \textit{Note:} The correct answer is: Carlisle
\item \textbf{Answer:} True
\\ \textit{Options:} A) True; B) False
\\ \textit{Note:} According to the passage: murder
\item \textbf{Answer:} False
\\ \textit{Options:} A) True; B) False
\\ \textit{Note:} The correct answer is: the world's first large-scale electrical supply network
\end{enumerate}

\section*{Long Form Response}
\begin{enumerate}[label=\arabic*.]
\setcounter{enumi}{10}
\item \textbf{Answer:} Roman times
\\ \textit{Note:} Expected answer: Roman times
\item \textbf{Answer:} Evie Hone
\\ \textit{Note:} Expected answer: Evie Hone
\end{enumerate}

\vfill
\noindent\textit{End of Answer Key}
\end{document}